%======================================================================
% 本科毕业设计（论文）模板（Ayaka Launcher 内置模板库）
% 引擎：XeLaTeX；编译命令：latexmk -xelatex main.tex
% 字体：Fandol（TeX Live 自带，无需安装系统字体）
% 结构：封面 → 摘要 → 目录 → 正文（章节文件在 sections/ 目录）→ 参考文献 → 致谢
% 说明：请将标注「TODO」的内容替换为实际信息；正文章节在 sections/ 下增删，
%       并同步修改 main.tex 中的 \input 行。
%======================================================================
\documentclass[zihao=-4,a4paper]{ctexart}

\usepackage{amsmath,amssymb}
\usepackage[a4paper,left=2.8cm,right=2.6cm,top=2.5cm,bottom=2.5cm]{geometry}
\usepackage{graphicx}
\usepackage{booktabs}
\usepackage{setspace}
\usepackage{hyperref}

\hypersetup{hidelinks, unicode, bookmarksnumbered}

%---------- 章节标题格式：第一章 / 1.1 ----------
\ctexset{
  section = {
    name        = {第,章},
    number      = \chinese{section},
    format      = \centering\zihao{3}\heiti,
    beforeskip  = 1.2ex plus 0.3ex,
    afterskip   = 1.2ex plus 0.3ex,
  },
  subsection = {
    format      = \zihao{4}\heiti,
    beforeskip  = 1.0ex plus 0.2ex,
    afterskip   = 0.8ex plus 0.2ex,
  },
}

%---------- 论文信息（TODO：按实际填写） ----------
\newcommand{\PaperTitle}{基于深度学习的图像识别方法研究}   % 论文题目
\newcommand{\SchoolName}{XX 大学}                          % 学校名称
\newcommand{\MajorName}{计算机科学与技术}                  % 专业名称
\newcommand{\StudentName}{张\quad 三}                      % 学生姓名
\newcommand{\StudentId}{2022010001}                        % 学号
\newcommand{\AdvisorName}{李四\quad 教授}                  % 指导教师
\newcommand{\FinishDate}{二〇二六年六月}                   % 完成日期

\begin{document}

%======================================================================
% 封面
%======================================================================
\begin{titlepage}
  \centering
  \vspace*{1.0cm}
  {\zihao{1}\heiti \SchoolName}\\[2.8cm]
  {\zihao{2}\heiti 本科毕业设计（论文）}\\[3.0cm]
  {\zihao{2}\heiti \PaperTitle}\\[3.5cm]
  {
    \zihao{3}
    \renewcommand{\arraystretch}{1.9}
    \begin{tabular}{r@{\quad}l}
      专\quad\quad 业： & \underline{\makebox[7.2cm][c]{\MajorName}} \\
      学\quad\quad 号： & \underline{\makebox[7.2cm][c]{\StudentId}} \\
      学生姓名：       & \underline{\makebox[7.2cm][c]{\StudentName}} \\
      指导教师：       & \underline{\makebox[7.2cm][c]{\AdvisorName}} \\
    \end{tabular}
  }
  \vspace{2.8cm}

  {\zihao{3}\FinishDate}
\end{titlepage}

%======================================================================
% 前置部分（摘要、目录）
%======================================================================
\pagenumbering{Roman}
%---------- 摘要 ----------
\begin{center}
  {\zihao{3}\heiti 摘\quad\quad 要}
\end{center}

\quad 随着人工智能技术的快速发展，深度学习在图像识别领域取得了显著成果。本文围绕
基于深度学习的图像识别方法展开研究。首先介绍了卷积神经网络的基本原理与典型结构；
其次针对传统方法在复杂场景下识别精度不足的问题，提出了一种结合注意力机制的改进
网络模型；最后在公开数据集上进行了对比实验，结果表明所提方法在识别准确率与推理
速度方面均优于基线模型。

\vspace{0.6em}
\noindent{\heiti 关键词：}深度学习；图像识别；卷积神经网络；注意力机制


\tableofcontents
\clearpage

%======================================================================
% 正文部分
%======================================================================
\pagenumbering{arabic}
%---------- 第一章 绪论 ----------
\section{绪\quad 论}

\subsection{研究背景与意义}

\quad 图像识别是计算机视觉领域的核心任务之一，广泛应用于人脸识别、医学影像分析、
工业质检、自动驾驶等场景。传统图像识别方法主要依赖人工设计的特征（如 SIFT、HOG）
与浅层分类器，特征表达能力有限，在复杂场景下的识别效果难以满足实际需求。

近年来，以卷积神经网络（Convolutional Neural Network, CNN）为代表的深度学习方法
通过端到端的特征自动学习，大幅提升了图像识别的精度，成为该领域的主流技术
路线\cite{ref-lecun2015}。研究基于深度学习的图像识别方法，对推动相关应用落地具有
重要的理论意义与实用价值。

\subsection{国内外研究现状}

\quad 自 2012 年 AlexNet 在 ImageNet 图像分类竞赛中取得突破以来，VGGNet、GoogLeNet、
ResNet\cite{ref-he2016} 等一系列深层网络结构相继被提出，网络深度与识别精度不断提升。
国内研究者在中文场景识别、轻量化网络设计等方向也取得了丰富成果\cite{ref-wang2020}。
然而，现有方法在目标尺度变化剧烈、背景干扰复杂的场景下仍存在识别精度不足的问题。

\subsection{本文主要工作}

\quad 本文的主要工作包括：
\begin{enumerate}
  \item 梳理卷积神经网络的基本原理与典型结构，分析现有图像识别方法的优势与不足；
  \item 提出一种结合注意力机制的改进网络模型，增强网络对关键区域的特征提取能力；
  \item 在公开数据集上开展对比实验，验证所提方法的有效性。
\end{enumerate}

\subsection{论文组织结构}

\quad 本文共分为四章：第一章为绪论，介绍研究背景、意义与国内外研究现状；第二章介绍
相关理论与技术基础；第三章详细阐述所提方法并给出实验结果与分析；第四章为总结与展望。

%---------- 第二章 相关理论与技术基础 ----------
\section{相关理论与技术基础}

\subsection{卷积神经网络}

\quad 卷积神经网络通过卷积层、池化层与全连接层的堆叠实现对图像特征的逐层提取。
设输入特征图为 $X \in \mathbb{R}^{H \times W \times C}$，卷积核为
$K \in \mathbb{R}^{k \times k \times C}$，则卷积层的输出可表示为
\begin{equation}
  Y_{i,j} = \sigma\left( \sum_{m=1}^{k} \sum_{n=1}^{k} \sum_{c=1}^{C}
    K_{m,n,c} \, X_{i+m-1,\,j+n-1,\,c} + b \right),
  \label{eq-conv}
\end{equation}
其中 $b$ 为偏置项，$\sigma(\cdot)$ 为非线性激活函数。式~\eqref{eq-conv} 表明，
卷积运算通过局部连接与权值共享有效降低了模型参数量。

\subsection{注意力机制}

\quad 注意力机制通过为特征图的不同空间位置分配不同权重，引导网络聚焦于对识别结果
贡献更大的区域。给定查询向量 $\boldsymbol{q}$ 与键值对集合
$\{(\boldsymbol{k}_i, \boldsymbol{v}_i)\}_{i=1}^{N}$，注意力权重定义为
\begin{equation}
  \alpha_i = \frac{\exp\left( \boldsymbol{q}^{\mathsf{T}} \boldsymbol{k}_i
    \middle/ \sqrt{d} \right)}
    {\sum_{j=1}^{N} \exp\left( \boldsymbol{q}^{\mathsf{T}} \boldsymbol{k}_j
    \middle/ \sqrt{d} \right)},
  \qquad
  \mathrm{Attn}(\boldsymbol{q}) = \sum_{i=1}^{N} \alpha_i \boldsymbol{v}_i,
\end{equation}
其中 $d$ 为缩放因子。常用的评价指标包括准确率、召回率与平均精度均值（mAP）。

\subsection{常用数据集}

\quad 本文实验所用公开数据集的基本情况如表~\ref{tab-datasets} 所示。

\begin{table}[htbp]
  \centering
  \caption{常用图像识别数据集概况}
  \label{tab-datasets}
  \begin{tabular}{cccc}
    \toprule
    数据集 & 图像数量 & 类别数 & 图像尺寸 \\
    \midrule
    MNIST     & 70\,000  & 10  & $28 \times 28$ \\
    CIFAR-10  & 60\,000  & 10  & $32 \times 32$ \\
    ImageNet  & $1.4 \times 10^{7}$ & 1\,000 & 不定 \\
    \bottomrule
  \end{tabular}
\end{table}

%---------- 第三章 方法与实验 ----------
\section{改进方法与实验分析}

\subsection{改进网络结构设计}

\quad 针对复杂场景下目标尺度变化剧烈、背景干扰严重导致的识别精度下降问题，本文在
基线网络的基础上引入通道—空间双重注意力模块：通道注意力用于筛选显著性特征通道，
空间注意力用于定位目标区域。改进后的网络在骨干网络的第 3、4 个残差阶段之后分别
插入注意力模块，并以跳跃连接保留原始特征信息，避免注意力加权造成的信息损失。

\subsection{实验设置}

\quad 实验环境为 Ubuntu 22.04，GPU 为 NVIDIA RTX 4090，深度学习框架为 PyTorch 2.2。
训练采用 SGD 优化器，初始学习率为 0.1，动量为 0.9，权重衰减为 $5 \times 10^{-4}$，
批大小为 128，共训练 200 个轮次，并在第 60、120、160 轮将学习率衰减为原来的
$1/10$。

\subsection{实验结果与分析}

\quad 本文方法与基线方法在 CIFAR-10 测试集上的对比结果如表~\ref{tab-results} 所示。
可以看出，引入注意力机制后，识别准确率提升了 1.6 个百分点，且参数量仅增加 3.2\%，
表明改进方法以极小的额外开销获得了明显的性能增益。

\begin{table}[htbp]
  \centering
  \caption{不同方法在 CIFAR-10 测试集上的对比结果}
  \label{tab-results}
  \begin{tabular}{lccc}
    \toprule
    方法 & 准确率（\%） & 参数量（M） & 推理时延（ms） \\
    \midrule
    ResNet-18（基线）        & 94.8 & 11.2 & 6.4 \\
    SE-ResNet-18             & 95.5 & 11.8 & 7.0 \\
    本文方法                 & 96.4 & 11.6 & 6.9 \\
    \bottomrule
  \end{tabular}
\end{table}

\quad 此外，消融实验表明：单独引入通道注意力可将准确率提升 0.9 个百分点，单独引入
空间注意力可提升 1.1 个百分点，二者结合可获得最优效果，说明通道与空间注意力在
特征增强中具有互补作用。

%---------- 第四章 总结与展望 ----------
\section{总结与展望}

\subsection{本文工作总结}

\quad 本文围绕基于深度学习的图像识别方法展开研究，完成了以下工作：
\begin{enumerate}
  \item 系统梳理了卷积神经网络与注意力机制的基本原理，分析了现有图像识别方法在
        复杂场景下精度不足的成因；
  \item 提出了一种结合通道—空间双重注意力机制的改进网络结构，兼顾识别精度与
        模型开销；
  \item 在公开数据集上完成了对比实验与消融实验，验证了所提方法的有效性。
\end{enumerate}

\subsection{不足与展望}

\quad 本文工作仍存在以下不足：其一，实验仅在 CIFAR-10 等中小规模数据集上开展，
方法在大规模数据集上的表现有待进一步验证；其二，注意力模块的插入位置与数量
依赖人工经验，后续可探索自动化神经架构搜索进行优化；其三，模型在边缘设备上的
部署与压缩尚需深入研究。未来将围绕上述方向继续完善所提方法，推动其在实际场景中
的落地应用。


%======================================================================
% 参考文献
%======================================================================
\begin{thebibliography}{9}
  \zihao{5}
  \bibitem{ref-lecun2015} LeCun Y, Bengio Y, Hinton G. Deep learning[J]. Nature, 2015, 521(7553): 436--444.
  \bibitem{ref-he2016} 何恺明, 张翔宇, 任少卿, 等. 深度残差学习中的图像识别[C]// 计算机视觉与模式识别会议. 拉斯维加斯: IEEE, 2016: 770--778.
  \bibitem{ref-wang2020} 王五, 赵六. 基于卷积神经网络的中文文本分类研究[J]. 计算机应用, 2020, 40(5): 1281--1287.
  \bibitem{ref-gb7714} 中华人民共和国国家标准化管理委员会. 信息与文献\ 参考文献著录规则: GB/T 7714—2015[S]. 北京: 中国标准出版社, 2015.
  \bibitem{ref-zhou2019} 周志华. 机器学习[M]. 北京: 清华大学出版社, 2019.
\end{thebibliography}

%======================================================================
% 致谢
%======================================================================
\section*{致\quad 谢}
\phantomsection
\addcontentsline{toc}{section}{致谢}
本论文的完成离不开许多人的帮助。首先感谢指导教师 \AdvisorName{} 老师在选题、
研究与写作过程中给予的悉心指导；感谢实验室同学在实验环境搭建与数据处理方面
提供的帮助；感谢家人和朋友一直以来的支持与鼓励。在此一并致以诚挚的谢意。

\end{document}
